% ===========================================================================
%  utente/07-marketplace/00-capitolo.tex
% ===========================================================================
\chapter{Marketplace}\label{cap:marketplace}

È il posto dove si comprano e si vendono strategie, approcci, overlay e
metodi di pesatura. È anche l'unica zona dell'applicazione che richiede un
account.

\figuraSchermata{1.0}{figure/schermate/marketplace.png}
  {Senza account. Il resto dell'applicazione continua a funzionare: la
   pagina lo dice.}
  {marketplace-senza-account}

\section{Perché qui serve un account e altrove no}\label{sez:perche-account}

Perché qui c'è qualcuno dall'altra parte. Costruire e provare strategie
avviene interamente sul proprio computer, e chiedere un account per farlo
non avrebbe senso. Comprare, vendere e dare una licenza richiede che il
sistema sappia chi è chi.

\section{Comprare}\label{sez:comprare}

\begin{pratica}{Dalla vetrina al sistema}
\begin{enumerate}
  \item Si apre la scheda dell'artefatto: descrizione, metriche dichiarate
        dal venditore, commenti, a volte un video.
  \item \pulsante{Buy} --- se è gratis l'accesso è immediato; se è a
        pagamento si apre il circuito di pagamento.
  \item Finito il pagamento si torna all'applicazione. Se l'accesso non
        risulta ancora attivo \textbf{non è un errore}: la conferma arriva
        un istante dopo. La pagina lo dice e invita ad aspettare.
  \item \pulsante{Install} --- l'artefatto viene scritto nel sistema e da
        quel momento compare negli elenchi come tutti gli altri.
\end{enumerate}
\end{pratica}

\begin{nota}{Prima di installare, i conflitti}
Se nel sistema esiste già qualcosa con lo stesso nome, l'applicazione lo
dice \emph{prima} di scrivere. Vale la pena leggerlo: un'installazione che
scavalca una strategia su cui si stava lavorando è un lavoro perso.
\end{nota}

\subsection{Se un acquisto non compare}

Esiste un comando apposta, nella scheda dell'artefatto, che spiega perché
non risulta acquistato. Distingue due cose che a prima vista sembrano
uguali: «non l'hai comprato» e «non sono riuscito a chiedere». Sono
opposte, e la seconda si risolve riprovando.

Se il pagamento è andato a buon fine ma l'accesso non si apre, c'è un
comando di riconciliazione: ricontrolla il pagamento presso il fornitore e,
solo se risulta incassato, completa l'operazione.

\section{Vendere}\label{sez:vendere}

Per pubblicare servono due cose, in quest'ordine:

\begin{enumerate}
  \item \textbf{Il marketplace deve essere aperto.} Durante la fase di
        prova è chiuso: nessuno carica niente. Gli articoli già in vetrina
        restano visibili.
  \item \textbf{Un abbonamento.} A marketplace aperto serve uno dei tre
        piani (\vediCap{account}).
\end{enumerate}

\begin{perche}{Perché un abbonamento anche per pubblicare gratis}
Non è per il denaro. Chi pubblica mette il proprio nome accanto a una
strategia che altri faranno girare sui propri soldi: un filtro all'ingresso
vale più di qualunque controllo a posteriori. Chi non è abbonato resta un
utente pieno --- sfoglia, scarica, compra e usa. Non può soltanto mettere
roba in vetrina.
\end{perche}

\noindent
Un artefatto pubblicato va online subito: non c'è una coda di approvazione.
La descrizione e il video si possono correggere dopo; il pacchetto di file,
no.

\section{Se un artefatto non rende come dichiarato}\label{sez:contestare}

Si apre una contestazione dalla scheda dell'artefatto. Non è una
discussione: il sistema esegue due verifiche automatiche.

\begin{center}
\begin{tabular}{@{}L{44mm}L{94mm}@{}}
\toprule
\textbf{Verifica} & \textbf{Che cosa confronta} \\
\midrule
\rigaTenue
\voce{Integrità} & I file installati con quelli venduti, uno per uno. \\
\voce{Performance} & Le metriche dichiarate dal venditore con quelle
  ottenute rieseguendo il calcolo sul proprio sistema. \\
\bottomrule
\end{tabular}
\end{center}

\noindent
La pratica è una conversazione fra acquirente, venditore e
amministrazione, con i risultati delle due verifiche allegati. Si può
ritirare in qualunque momento.

\begin{nota}{Due risultati diversi non significano un inganno}
Una strategia può rendere diversamente perché i dati sono diversi: un
universo di strumenti diverso, un periodo diverso, costi diversi. La
verifica serve a stabilire \emph{di quanto} differiscono, che è il punto di
partenza di qualunque discussione sensata.
\end{nota}
